\documentclass[11pt]{article}
\usepackage[margin=1in]{geometry}
\usepackage{amsmath,amssymb,amsthm}
\usepackage{graphicx}
\usepackage[hidelinks]{hyperref}

\newtheorem{theorem}{Theorem}
\newtheorem{lemma}{Lemma}
\newtheorem{remark}{Remark}

\newcommand{\K}{\mathbb K}
\newcommand{\cA}{\mathcal A}
\newcommand{\cL}{\mathcal L}
\newcommand{\liminfn}{\underline{\lim}}

\title{A wild-dynamics operator on every infinite-dimensional Banach space}
\author{fa\_banach\_001, lane 7}
\date{2026-06-20}

\begin{document}
\maketitle

\begin{abstract}
Aug\'e proved that every infinite-dimensional separable Banach space admits an
operator \(R=I+K\), with \(K\) compact and even nuclear, for which
\[
  A_R=\{x:\|R^t x\|\to\infty\}
\]
has non-empty interior but is not dense; equivalently, \(A_R\) and a recurrent
set \(B_R\) form a partition of the space.  Remark 3.8 of the source paper
conjectures that the same construction should be possible in arbitrary Banach
spaces.  This packet proves that conjectural extension.  The only extra
ingredient is to replace the separable bounded biorthogonal system by a
normalized basic sequence and to choose the periods so fast that the periodic
part returns to the identity in operator norm.
\end{abstract}

\section{Source question}

The source is Jean-Matthieu Aug\'e, \emph{Linear operators with wild dynamics},
arXiv:1204.2044.  Theorem 1.1 proves the desired wild-dynamics construction
for separable Banach spaces.  The open-ended extension appears in Remark 3.8
on page 8:

\begin{center}
\includegraphics[width=\linewidth]{figures/open_problem_crop.png}
\end{center}

In the notation of the source paper, \(\cA_X\) denotes the class of operators
for which \(A_R\) and
\[
  B_R=\{x:\liminfn_{t\to\infty}\|R^t x-x\|=0\}
\]
both have non-empty interior and form a partition of \(X\).

\section{Proof intuition}

The separable proof uses an Ovsepian--Pelczynski bounded biorthogonal sequence
whose linear span is dense in \(X\).  Density is used to show that a diagonal
operator \(S\), with eigenvalues roots of unity tending rapidly to \(1\),
satisfies \(S^{2m_n}x\to x\) for every \(x\in X\).

For a nonseparable space we cannot ask a countable sequence to span densely.
However, every infinite-dimensional Banach space contains a normalized basic
sequence \((e_n)\), and its coordinate functionals extend to uniformly bounded
functionals on \(X\).  If the periods \(m_n\) are chosen much faster than in
the separable proof, then \(S^{2m_n}\to I\) in operator norm, not merely
strongly on a dense span.  That operator-norm return is exactly what removes
the separability hypothesis.  The growth side of Aug\'e's proof is unchanged:
a finite-dimensional separating lemma supplies functionals \(f_n\) which grow
to infinity off a prescribed union of lines and return arbitrarily close to
zero on it.

\section{Preliminaries}

We use two standard facts.  First, every infinite-dimensional Banach space
contains a normalized basic sequence.  Hence there are vectors
\((e_n)\subset X\) with \(\|e_n\|=1\) and functionals \(e_n^*\in X^*\) such
that
\[
  e_n^*(e_m)=\delta_{nm},
  \qquad C:=\sup_n\|e_n^*\|<\infty .
\]
Indeed the coordinate functionals of the basic sequence are uniformly bounded
on its closed span, and Hahn--Banach extends them to \(X\).

Second, we use Aug\'e's finite-dimensional separating lemma
\cite[Proposition 2.1]{Auge}.  In the two-dimensional space
\(E=\operatorname{span}\{e_1,e_2\}\), let
\[
  F=\{a e_1+b e_2: |b|\le |a|\}.
\]
This is a closed proper union of one-dimensional \(\K\)-subspaces of \(E\),
and both \(F\) and \(E\setminus F\) have non-empty relative interior.  The
lemma gives linear forms \(f_k\in E^*\), indexed below by \(k\ge 3\), such
that
\[
  y\notin F \Rightarrow |f_k(y)|\to\infty,
  \qquad
  y\in F \Rightarrow \liminfn_{k\to\infty}|f_k(y)|=0 .
\]

Choose positive integers \(m_k\), with \(m_1=m_2=1\) and \(m_k\mid m_{k+1}\),
so rapidly increasing that
\begin{equation}\label{eq:growth}
  \sum_{k=3}^\infty \frac{\|f_k\|}{m_{k-1}}<\infty,\qquad
  \sum_{k=3}^\infty \frac{m_{k-2}}{m_{k-1}}\|f_k\|<\infty,\qquad
  \alpha_n:=\sum_{k>n}\frac{m_n}{m_k}\longrightarrow 0 .
\end{equation}
Such a sequence is obtained recursively, choosing each next multiple large
enough to make the displayed tails as small as desired.

\section{The complex case}

\begin{theorem}\label{thm:complex}
Let \(X\) be an infinite-dimensional complex Banach space.  There is
\(R\in\cL(X)\) such that \(R=I+N\) with \(N\) nuclear and
\[
  A_R=\{x:Px\notin F\},\qquad B_R=\{x:Px\in F\},
\]
where \(P x=e_1^*(x)e_1+e_2^*(x)e_2\).  In particular \(A_R\) and \(B_R\)
have non-empty interior and form a partition of \(X\).
\end{theorem}

\begin{proof}
For \(k\ge 3\) put \(\lambda_k=\exp(\pi i/m_k)\), and define
\[
  Sx=x+\sum_{k=3}^\infty (\lambda_k-1)e_k^*(x)e_k,
\]
\[
  Rx=Sx+\sum_{k=3}^\infty \frac{1}{m_{k-1}}f_k(Px)e_k .
\]
The two series converge in operator norm by \eqref{eq:growth} and
\(|\lambda_k-1|\le \pi/m_k\).  Moreover
\[
  R-I=\sum_{k=3}^\infty
  e_k\otimes\left((\lambda_k-1)e_k^*+\frac{1}{m_{k-1}}f_k\circ P\right),
\]
and the sum of the norms of these rank-one summands is finite.  Thus
\(R-I\) is nuclear.

For \(t\ge 1\), set
\[
  \lambda_{k,t}=1+\lambda_k+\cdots+\lambda_k^{t-1}.
\]
Since \(Pe_k=0\) for \(k\ge 3\) and \(PS^t=P\), induction gives
\begin{equation}\label{eq:iterate}
  R^t x=S^t x+\sum_{k=3}^\infty
        \frac{\lambda_{k,t}}{m_{k-1}}f_k(Px)e_k .
\end{equation}
Also
\[
  S^t x=x+\sum_{k=3}^\infty(\lambda_k^t-1)e_k^*(x)e_k .
\]
If \(q_n=2m_n\), then \(\lambda_k^{q_n}=1\) for \(k\le n\), and hence
\[
  \|S^{q_n}-I\|
  \le C\sum_{k>n}|\lambda_k^{q_n}-1|
  \le 2\pi C\sum_{k>n}\frac{m_n}{m_k}
  =2\pi C\alpha_n\to0 .
\]

Let \(y=Px\notin F\).  For \(m_{k-1}\le t\le m_k\),
\[
  |\lambda_{k,t}|
  =\frac{|\sin(\pi t/(2m_k))|}{|\sin(\pi/(2m_k))|}
  \ge \frac{2}{\pi}m_{k-1}.
\]
Applying \(e_k^*\) to \eqref{eq:iterate},
\[
  C\|R^t x\|
  \ge |e_k^*(R^t x)|
  \ge \frac{|\lambda_{k,t}|}{m_{k-1}}|f_k(y)|-|e_k^*(x)|
  \ge \frac{2}{\pi}|f_k(y)|-C\|x\|.
\]
As \(t\to\infty\), the corresponding \(k\to\infty\), and
\(|f_k(y)|\to\infty\).  Thus \(\|R^t x\|\to\infty\), so \(x\in A_R\).

Now let \(y=Px\in F\).  Choose \(k_j\to\infty\) such that
\(|f_{k_j}(y)|\to0\), and put \(q_j=2m_{k_j-1}\).  The \(S\)-part satisfies
\(S^{q_j}x\to x\).  For the perturbation part in \eqref{eq:iterate}, the
terms with \(k<k_j\) vanish because \(\lambda_{k,q_j}=0\).  The \(k=k_j\)
term has norm at most \(2|f_{k_j}(y)|\).  The tail \(k>k_j\) is bounded by
\[
  2\|P\|\|x\|\sum_{k>k_j}\frac{m_{k_j-1}}{m_{k-1}}\|f_k\|
  \le
  2\|P\|\|x\|\sum_{k>k_j}\frac{m_{k-2}}{m_{k-1}}\|f_k\|,
\]
which tends to \(0\) by \eqref{eq:growth}.  Therefore
\(\liminfn_t\|R^t x-x\|=0\), so \(x\in B_R\).

We have proved \(P^{-1}(E\setminus F)\subset A_R\) and
\(P^{-1}(F)\subset B_R\).  These two preimages partition \(X\), and a vector
in \(B_R\) cannot have norm tending to infinity along the whole orbit.  Hence
the displayed equalities hold.  Since \(P\) is a continuous projection onto
the two-dimensional space \(E\), and both \(F\) and \(E\setminus F\) have
non-empty relative interior, both \(A_R\) and \(B_R\) have non-empty interior
in \(X\).
\end{proof}

\section{The real case}

\begin{theorem}\label{thm:real}
The same conclusion holds for every infinite-dimensional real Banach space.
\end{theorem}

\begin{proof}
Use the same first two vectors \(e_1,e_2\), the same projection \(P\), the
same set \(F\subset \operatorname{span}\{e_1,e_2\}\), and separating forms
\((f_k)_{k\ge2}\).  Group the remaining basic vectors into pairs
\((e_{2k-1},e_{2k})\), \(k\ge2\), and choose the periods \(m_k\) satisfying
the analogues of \eqref{eq:growth}.

Let \(\theta_k=\pi/m_k\).  Define \(S\) so that it is the identity outside the
chosen coordinates and, on each pair, acts by the rotation
\[
  e_{2k-1}\mapsto \cos\theta_k\,e_{2k-1}+\sin\theta_k\,e_{2k},
  \qquad
  e_{2k}\mapsto -\sin\theta_k\,e_{2k-1}+\cos\theta_k\,e_{2k}.
\]
In operator form this differs from the identity by an absolutely summable
rank-one series, since \(\sum_k |\sin\theta_k|<\infty\) and
\(\sum_k |\cos\theta_k-1|<\infty\).  Set
\[
  Rx=Sx+\sum_{k=2}^\infty \frac{1}{m_{k-1}}f_k(Px)e_{2k}.
\]
Again \(R-I\) is nuclear.

Let
\[
  \mu_{k,t}=\sum_{\ell=0}^{t-1}S^\ell e_{2k}.
\]
The same induction as before gives
\[
  R^t x=S^t x+\sum_{k=2}^\infty \frac{1}{m_{k-1}}f_k(Px)\mu_{k,t}.
\]
The rotation period is \(2m_k\), so \(\mu_{k,2m_n}=0\) whenever \(k\le n\).
The choice of periods also gives \(S^{2m_n}\to I\) in operator norm: for
\(k>n\), the norm of the difference between the \(2m_n\)-th rotation and the
identity is \(O(m_n/m_k)\), and the tail is controlled by \(\alpha_n\to0\).

For \(m_{k-1}\le t\le m_k\), the two coordinates of \(\mu_{k,t}\) are the real
and imaginary parts, up to signs, of
\(\sum_{\ell=0}^{t-1}e^{i\ell\theta_k}\).  Hence at least one coordinate has
absolute value at least
\[
  \frac{1}{\sqrt2}
  \frac{|\sin(\pi t/(2m_k))|}{|\sin(\pi/(2m_k))|}
  \ge \frac{2}{\pi\sqrt2}m_{k-1}.
\]
Applying the corresponding coordinate functional \(e_{2k-1}^*\) or
\(e_{2k}^*\) yields
\[
  C\|R^t x\|\ge c\,|f_k(Px)|-C\|x\|
\]
with an absolute constant \(c>0\).  This proves divergence for \(Px\notin F\).

If \(Px\in F\), choose \(k_j\) with \(f_{k_j}(Px)\to0\) and take
\(q_j=2m_{k_j-1}\).  The blocks with \(k<k_j\) vanish at time \(q_j\), the
\(k_j\)-block is bounded by a constant times \(|f_{k_j}(Px)|\), and the tail
is bounded by
\[
  C'\|P\|\|x\|\sum_{k>k_j}\frac{m_{k-2}}{m_{k-1}}\|f_k\|,
\]
which tends to zero.  Thus \(R^{q_j}x\to x\).  The same partition and interior
argument as in the complex case completes the proof.
\end{proof}

\section{Verification and novelty check}

The proof uses no numerical computation.  The packet includes
\texttt{code/crop\_open\_problem.py}, which only renders the page-8 source
crop from \texttt{source\_paper.pdf}.

Local duplicate checks in the run indexes found no previous packet for
arXiv:1204.2044 or for this arbitrary-Banach-space strengthening.  External
web searches on 2026-06-20 used the exact title, arXiv id, Remark 3.8 wording,
``operators with the properties stated in Theorem 1.1 arbitrary Banach
spaces'', ``Prajitura's conjecture arbitrary Banach spaces'', and nearby
phrases involving \(A_R\), \(I+K\), and nuclear perturbations.  These searches
found the source paper and the earlier Hajek--Smith symmetric-basis paper, but
no later exact answer to Remark 3.8.  This is a bounded novelty check, not a
claim of exhaustive bibliographic certainty.

\section{Human review recommendation}

Review the operator-norm return estimate \(S^{2m_n}\to I\), since that is the
new point replacing separability.  Also check that the real rotation-block
translation preserves the lower-coordinate estimate and the nuclearity of
\(R-I\).  If those estimates are accepted, the packet gives a full positive
answer to Aug\'e's arbitrary Banach space conjecture in Remark 3.8.

\begin{thebibliography}{9}

\bibitem{Auge}
J.-M. Aug\'e,
\emph{Linear operators with wild dynamics},
arXiv:1204.2044, 2012.

\bibitem{HajekSmith}
P. H\'ajek and R. J. Smith,
\emph{Operator machines on directed graphs},
Integral Equations Operator Theory 67 (2010), 15--31.

\bibitem{Singer}
I. Singer,
\emph{Bases in Banach Spaces I},
Springer, 1970.

\end{thebibliography}

\end{document}
